\chapter{Scale, Variety and Firms}\label{ch:b3:scale-variety-firms}

The two previous chapters explained trade by differences between countries: in
technology, in endowments. Yet much of world trade is between similar rich
countries, and much of it within the same industries: Germany sells cars to
France and buys cars from it. And inside every industry, most firms do not trade
at all. This chapter explains both facts with two models. In the first, every
firm makes its own variety, consumers like variety, and production has fixed
costs, so trade lets each country make fewer varieties at larger scale while its
consumers enjoy more. In the second, firms differ in productivity, and only the
best can pay the cost of selling abroad. The opening chart shows how few they are.

\begin{openingchart}{Share of firms that export, by manufacturing industry, United
States, 2002 (percent). The dotted line is all manufacturing, 18~\%.}{Bernard,
Jensen, Redding and Schott (2007), Table~2}
\begin{tikzpicture}
\begin{axis}[omrecord, xbar, bar width=5pt, xmin=0, xmax=45, ymin=0.3, ymax=21.7,
    y dir=reverse, ytick=data,
    yticklabels from table={figdata/bachelor-3/09-scale-variety-firms/out/exporters.csv}{industry},
    y tick label style={font=\footnotesize}, xlabel={Percent of the industry's firms that export},
    grid=none, xmajorgrids, width=0.62\linewidth, height=0.72\linewidth,
    nodes near coords, every node near coord/.append style={font=\scriptsize}]
  \addplot[fill=omDef!70, draw=omDef] table[x=exporting, y=pos]
    {figdata/bachelor-3/09-scale-variety-firms/out/exporters.csv};
  \draw[guide] (axis cs:18,0.3) -- (axis cs:18,21.7);
\end{axis}
\end{tikzpicture}
\end{openingchart}

% ledger: b3sv:bjrs.t2.all, b3sv:bjrs.t2.334, b3sv:bjrs.t2.339, b3sv:bjrs.t2.323
Fewer than one United States manufacturing firm in five exported in 2002. The share
ranged from 38~\% in computers and electronics and in electrical equipment to
5~\% in printing and 2~\% in miscellaneous manufacturing --- but every industry
had exporters, and in every industry most firms sold only at home. And the firms
that did export sold most of their output at home too: their exports were, on
average, 14~\% of their shipments.

\begin{recall}
The Year 1 volume taught monopoly pricing, the Lerner index $(p - c)/p = 1/|\varepsilon|$,
economies of scale and free entry; the grades 10--12 volume introduced
\omterm{def:b3:scale-variety-firms:monopolistic-competition}{monopolistic competition}. Chapters \ref{ch:b3:ricardian-trade} and
\ref{ch:b3:factor-endowments} built the trade models this chapter adds to.
\end{recall}

\begin{notation}[Symbols of this chapter]\label{not:b3:scale-variety-firms:symbols}
$\sigma > 1$ is the elasticity of substitution between varieties, $n$ the number
of varieties (and of firms), $P$ the price index of the industry and $E$ spending
on it. A firm has fixed cost $F$ and marginal cost $c$, both in labour, with the
wage set to 1. $\tau \ge 1$ units must be shipped for one to arrive (\omterm{def:b3:ricardian-trade:iceberg}{iceberg
trade cost}); $\lambda = \tau^{1-\sigma}$ measures the freeness of trade (not an import
share). In the last sections $\varphi$ is a firm's productivity, $\varphi^*$ the
domestic cut-off and $\varphi_x$ the \omterm{def:b3:scale-variety-firms:melitz}{export cut-off}.
\end{notation}

\section{Trade inside industries}\label{sec:b3:scale-variety-firms:intra}

\begin{definition}[Intra-industry trade; Grubel--Lloyd index]\label{def:b3:scale-variety-firms:intra-industry}
\index{intra-industry trade}\index{Grubel--Lloyd index}
\emph{Intra-industry trade} is a country's exports and imports of goods of the
same industry. For an industry with exports $X$ and imports $M$, the
\emph{Grubel--Lloyd index} is $\mathrm{GL} = 1 - |X - M|/(X + M)$.
\end{definition}

\begin{measure}[The Grubel--Lloyd index]
The index is 0 when trade in the industry goes one way only, as Ricardo and
Heckscher--Ohlin predict, and 1 when exports equal imports. A country's
aggregate index weights the industries by their trade: $1 - \sum_i |X_i - M_i| /
\sum_i (X_i + M_i)$. Its value depends on how finely industries are drawn: at a
coarse level (``machinery''), a country exports some kinds and imports others,
and the index is high; at a fine one (``ball bearings of a given size''),
part of that trade becomes inter-industry and the index falls. A high index
between countries with similar endowments is the fact that endowment models
cannot produce (\cref{exo:b3:scale-variety-firms:2}).
\end{measure}

\section{Love of variety and monopolistic competition}\label{sec:b3:scale-variety-firms:variety}

\begin{definition}[Love of variety]\label{def:b3:scale-variety-firms:variety}
\index{love of variety}\index{CES preferences}
Consumers have a \emph{love of variety} when, for given spending, they are
better off with more varieties of a good. The workhorse is \emph{CES
preferences} over $n$ varieties,
\begin{equation}\label{eq:b3:sv:ces}
  U = \Bigl(\sum_{i=1}^n x_i^{(\sigma - 1)/\sigma}\Bigr)^{\sigma/(\sigma - 1)}, \qquad \sigma > 1 .
\end{equation}
\end{definition}

Maximising \eqref{eq:b3:sv:ces} under the budget $\sum_i p_i x_i = E$ gives the
demand for variety $i$ (\cref{exo:b3:scale-variety-firms:10}),
\begin{equation}\label{eq:b3:sv:demand}
  x_i = \frac{p_i^{-\sigma}}{P^{1-\sigma}}\, E, \qquad P = \Bigl(\sum_j p_j^{1-\sigma}\Bigr)^{1/(1-\sigma)} ,
\end{equation}
and utility $U = E/P$. With $n$ varieties at the same price $p$, $P = n^{1/(1-\sigma)}p$,
so $U = n^{1/(\sigma-1)} E/p$: at a given price and spending, more varieties mean
more utility. The \omterm{def:b3:scale-variety-firms:variety}{love of variety} is stronger the smaller $\sigma$, the less
alike the varieties.

\begin{definition}[Monopolistic competition]\label{def:b3:scale-variety-firms:monopolistic-competition}
\index{monopolistic competition}
\emph{Monopolistic competition} is a market in which each firm is the only
seller of its variety and so sets its price, but free entry drives profits to
zero, and each firm is too small for its price to move the industry's price
index $P$.
\end{definition}

\begin{proposition}[A constant markup]\label{prop:b3:scale-variety-firms:markup}
Under CES demand and \omterm{def:b3:scale-variety-firms:monopolistic-competition}{monopolistic competition}, every firm charges $p = \dfrac{\sigma}{\sigma - 1}\, c$.
\end{proposition}

\begin{proof}
By \eqref{eq:b3:sv:demand}, with $P$ and $E$ taken as given, the firm faces demand
$x = A p^{-\sigma}$ with $A$ fixed: its price elasticity is $-\sigma$. Profit
$(p - c)Ap^{-\sigma}$ is maximised where $Ap^{-\sigma} - \sigma(p - c)Ap^{-\sigma - 1} = 0$,
that is $p = \sigma(p - c)$, or $p = \sigma c/(\sigma - 1)$: the Lerner condition $(p -
c)/p = 1/\sigma$.
\end{proof}

\section{The Krugman model of trade}\label{sec:b3:scale-variety-firms:krugman}

Let a country have $L$ workers. Each firm needs $F$ workers to start and $c$
workers per unit; it charges the markup price of
\cref{prop:b3:scale-variety-firms:markup}.

\begin{proposition}[Firm size, variety and the gains from trade]\label{prop:b3:scale-variety-firms:krugman}
With free entry, every firm makes $\bar x = (\sigma - 1)F/c$ units and employs
$\sigma F$ workers, and there are $n = L/(\sigma F)$ firms. Free trade between two
such countries, with $L$ and $L^*$ workers, leaves each firm's size and price
unchanged, raises the number of varieties each consumer buys to $n + n^* = (L +
L^*)/(\sigma F)$, and raises every worker's real wage by the factor $\bigl((L +
L^*)/L\bigr)^{1/(\sigma-1)}$.
\end{proposition}

\begin{proof}
Zero profit: $(p - c)\bar x = F$ with $p - c = c/(\sigma - 1)$, so $\bar x = (\sigma -
1)F/c$, and the firm's employment is $F + c\bar x = \sigma F$. Full employment gives
$n\sigma F = L$. Neither condition involves the size of the market, so opening to
trade changes neither $\bar x$ nor $p$, only the number of firms the consumer can
buy from: $n + n^*$. With the wage at 1, the real wage is $U$ per unit of spending,
$1/P = n^{1/(\sigma-1)}/p$; it rises with the number of varieties to the power
$1/(\sigma - 1)$.
\end{proof}

Similar countries trade, and trade within industries, because each makes
different varieties; each country's industry shrinks in number of firms
relative to the combined market, and both gain. In this version firm size is
fixed by the CES markup; with demand that becomes more elastic as varieties
multiply, trade also makes firms larger and prices lower.

\begin{record}[Gains from new varieties]
Broda and Weinstein counted the varieties --- a product from a given country ---
that the United States imported: their number rose by a factor of four from 1972
to 2001. Because the official import price index ignores new varieties, it
overstated import price inflation by about 1.2~\% a year, and the welfare gain
from the growth in imported varieties alone came to 2.8~\% of GDP (the working-paper
estimate). % ledger: b3sv:bw.variety
\end{record}

\section{The home-market effect}\label{sec:b3:scale-variety-firms:home-market}

\begin{definition}[Home-market effect]\label{def:b3:scale-variety-firms:home-market}
\index{home-market effect}
The \emph{home-market effect} is the prediction that, with increasing returns and
trade costs, a country's share of an industry's production rises more than in
proportion to its share of the industry's demand, so that the larger market
becomes a net exporter of the industry's goods.
\end{definition}

Let Home and Foreign spend $E$ and $E^*$ on a differentiated good, and let an
outside good, traded freely, keep both wages at 1, so that all varieties sell at
the same factory price $p$. Shipping costs $\tau$, so a foreign consumer pays $\tau
p$. A Home firm's sales are, by \eqref{eq:b3:sv:demand},
\[
  R = \frac{E}{n + \lambda n^*} + \frac{\lambda E^*}{\lambda n + n^*}, \qquad \lambda = \tau^{1-\sigma},
\]
and a Foreign firm's $R^* = \lambda E/(n + \lambda n^*) + E^*/(\lambda n + n^*)$. Free entry
requires each to equal $\bar R = \sigma F$.

\begin{proposition}[The home-market effect]\label{prop:b3:scale-variety-firms:hme}
If both countries make the good, Home's share of the firms is
\[
  s_n = \frac{n}{n + n^*} = \frac12 + \frac{1 + \lambda}{1 - \lambda}\Bigl(s_E - \frac12\Bigr),
  \qquad s_E = \frac{E}{E + E^*},
\]
so a one-point rise in Home's share of demand raises its share of production by
$(1 + \lambda)/(1 - \lambda) > 1$ points, and the country with the larger market is a net
exporter of the good.
\end{proposition}

\begin{proof}
Write $a = E/(n + \lambda n^*)$ and $b = E^*/(\lambda n + n^*)$. Free entry gives $a + \lambda b =
\lambda a + b = \bar R$; subtracting, $(1 - \lambda)(a - b) = 0$, so $a = b = \bar R/(1 + \lambda)$.
Hence $n + \lambda n^* = E(1 + \lambda)/\bar R$ and $\lambda n + n^* = E^*(1 + \lambda)/\bar R$. Solving,
$n = (E - \lambda E^*)/\bigl(\bar R(1 - \lambda)\bigr)$, $n^* = (E^* - \lambda E)/\bigl(\bar R(1 -
\lambda)\bigr)$, and $n + n^* = (E + E^*)/\bar R$. Dividing, $s_n = (s_E(1 + \lambda) - \lambda)/(1
- \lambda)$, which rearranges to the formula. With $s_E > \frac12$, $s_n > s_E$: Home's firms
sell $s_n(E + E^*)$, more than its spending $s_E(E + E^*)$, so it exports on net.
\end{proof}

\begin{omchart}{Home's share of the industry's firms against its share of
world spending on the industry, in the two-country Krugman model, for three
levels of the freeness of trade $\lambda$.}{Model of this section; our computation}
\begin{tikzpicture}
\begin{axis}[omrecord, width=0.55\linewidth, height=0.5\linewidth, xmin=0, xmax=1,
  ymin=0, ymax=1, xlabel={Home's share of spending $s_E$}, ylabel={Home's share of firms $s_n$},
  legend pos=north west, legend style={font=\footnotesize}]
  \addplot[omThm, thick, no marks, sharp plot] table[x=s_e, y=lam2] {figdata/bachelor-3/09-scale-variety-firms/out/hme.csv};
  \addplot[omDef, thick, no marks, sharp plot] table[x=s_e, y=lam5] {figdata/bachelor-3/09-scale-variety-firms/out/hme.csv};
  \addplot[black, thick, no marks, sharp plot] table[x=s_e, y=lam8] {figdata/bachelor-3/09-scale-variety-firms/out/hme.csv};
  \addplot[guide, sharp plot] table[x=s_e, y=s_e] {figdata/bachelor-3/09-scale-variety-firms/out/hme.csv};
  \legend{$\lambda = 0.2$, $\lambda = 0.5$, $\lambda = 0.8$}
\end{axis}
\end{tikzpicture}
\end{omchart}

The dotted diagonal is production in proportion to demand. The steeper the line,
the stronger the effect, and it grows as trade gets cheaper: when shipping is
nearly free, a small edge in market size draws the whole industry into one
country. The prediction separates this model from the endowment models, in which
a country with a strong taste for a good imports it
(\cref{exo:b3:scale-variety-firms:4}).

\section{Firms that differ: the Melitz model}\label{sec:b3:scale-variety-firms:melitz}

In the Krugman model every firm exports. The opening chart says most do not.

\begin{definition}[Heterogeneous-firms model; export cut-off]\label{def:b3:scale-variety-firms:melitz}
\index{heterogeneous-firms model}\index{export cut-off}
In the \emph{heterogeneous-firms model} (Melitz), an entrant pays a sunk entry
cost $f_e$, then draws a productivity $\varphi$, its output per worker. To produce
it pays a fixed cost $f$ each period; to export, a further fixed cost $f_x$ and the
trade cost $\tau$. Demand is CES and entry is free. Firms with $\varphi$ below the
\emph{domestic cut-off} $\varphi^*$ exit; firms with $\varphi$ above the \emph{export
cut-off} $\varphi_x$ export as well.
\end{definition}

A firm's revenue is proportional to $\varphi^{\sigma-1}$, and its profit is revenue
divided by $\sigma$, less fixed costs. At the domestic cut-off, profit is zero, so
revenue is $\sigma f$, and a firm with $\varphi \ge \varphi^*$ has domestic profit $f\bigl((\varphi/\varphi^*)^{\sigma-1}
- 1\bigr)$. Its foreign revenue is $\lambda$ times its domestic revenue (in a market
of the same size), so it exports when $\lambda\, \sigma f(\varphi/\varphi^*)^{\sigma-1}/\sigma
\ge f_x$:
\begin{equation}\label{eq:b3:sv:export-cutoff}
  \varphi_x = \tau \Bigl(\frac{f_x}{f}\Bigr)^{1/(\sigma-1)} \varphi^* .
\end{equation}
Free entry sets the expected profit of an entrant equal to $f_e$. When
productivity is Pareto-distributed, $\Pr(\varphi > z) = z^{-k}$ for $z \ge 1$ with $k > \sigma
- 1$, the expected profit can be integrated in closed form (as in the weekend
problem), and the cut-off follows.

\begin{proposition}[Opening raises the cut-off]\label{prop:b3:scale-variety-firms:cutoff}
In the \omterm{def:b3:scale-variety-firms:melitz}{heterogeneous-firms model} with Pareto productivity, the closed-economy
cut-off satisfies $\varphi^{*k} = f(\sigma - 1)/\bigl((k - \sigma + 1)f_e\bigr)$; opening to
trade with a symmetric country multiplies $\varphi^{*k}$ by $1 + \tau^{-k}(f/f_x)^{k/(\sigma-1) -
1}$. The least productive firms exit, the most productive export, and average
productivity rises.
\end{proposition}
\admitted

The listing simulates 20{,}000 entrants in the economy of the weekend problem:
$\sigma = 3$, $k = 4$, $f = 40$, $f_e = 2.5$, $\tau = 1.2$ and $f_x = 90$. % figdata: furn_params

\omcode{code/bachelor-3/09-scale-variety-firms/listings/melitz_selection.py}{1}{26}{Selection in the Melitz model, closed and open.}\label{lst:b3:scale-variety-firms:melitz}

\omoutput{code/bachelor-3/09-scale-variety-firms/outputs/melitz_selection.txt}{Opening pushes the weakest out.}\label{out:b3:scale-variety-firms:melitz}

In autarky one entrant in sixteen produces. With trade, the cut-off rises from 2
to 2.10: about 18~\% of the firms that produced in autarky exit, and in this
sample 8.7~\% of the producers export (the formula's population share is 9.5~\%).
% output: melitz_selection
% figdata: furn_results
No firm has become more productive, yet the average producer has: the least
productive firms are gone. The figure below places the three cut-offs on the
productivity line, and the chart shows what happens to each firm's sales.

\begin{omfigure}
\begin{tikzpicture}[econ, x=1.6cm, y=1cm]
  \draw[econ axes] (1,0) -- (6.5,0) node[right] {$\varphi$};
  \foreach \x/\l in {1/1, 2/2, 3/3, 4/4, 5/5, 6/6} \draw (\x,0.08) -- (\x,-0.08) node[below, font=\scriptsize] {\l};
  \draw[omThm, thick] (2,0) -- (2,0.9) node[above, font=\scriptsize] {$\varphi^*$ closed};
  \draw[omDef, thick] (2.0995,0) -- (2.0995,0.5) node[right, font=\scriptsize] {$\varphi^*$ open};
  \draw[black, thick] (3.7791,0) -- (3.7791,0.9) node[above, font=\scriptsize] {$\varphi_x$};
  \node[font=\scriptsize] at (1.5,-0.75) {never produce};
  \node[font=\scriptsize] at (2.95,-0.75) {sell at home only};
  \node[font=\scriptsize] at (5.1,-0.75) {sell at home and abroad};
\end{tikzpicture}
\par\smallskip{\small The productivity line in the economy of the listing. Firms
just above the closed-economy cut-off exit when the country opens.}
\end{omfigure}

\begin{omchart}{Revenue of each simulated firm against its productivity, before
and after opening to trade (thousands of dalens a year, log scale; the producers among 10{,}000 entrants,
up to a productivity of 6).}{Simulated: Melitz model of the weekend problem}
\begin{tikzpicture}
\begin{axis}[omscatter, width=0.62\linewidth, height=0.45\linewidth, ymode=log,
  log ticks with fixed point, xmin=1.9, xmax=6.1, ymin=80, ymax=2000,
  xlabel={Productivity $\varphi$}, ylabel={Revenue}, legend pos=north west,
  legend style={font=\footnotesize}]
  \addplot[only marks, mark=o, mark size=1.6pt, omThm] table[x=phi, y=closed] {figdata/bachelor-3/09-scale-variety-firms/out/furn_firms.csv};
  \addplot[only marks, mark=*, mark size=1.2pt, omDef] table[x=phi, y=open] {figdata/bachelor-3/09-scale-variety-firms/out/furn_firms.csv};
  \draw[guide] (axis cs:3.7791,80) -- (axis cs:3.7791,2000);
  \legend{closed, open}
\end{axis}
\end{tikzpicture}
\end{omchart}

Firms just above the old cut-off have no filled dot: they exit. Firms between the
two cut-offs sell less than before, because imported varieties take part of
their market; firms beyond the \omterm{def:b3:scale-variety-firms:melitz}{export cut-off} (the dotted line) jump up, by a
factor $1 + \lambda = 1.69$ at given productivity. Market share moves from the least
to the most productive firms, and that move, not any firm's learning, raises the
industry's productivity (\cref{exo:b3:scale-variety-firms:11}).

\begin{model}{Monopolistic competition and trade}
\assumes CES \omterm{def:b3:scale-variety-firms:variety}{love of variety}; increasing returns from a fixed cost; free entry;
identical firms.
\predicts Similar countries trade, and trade within industries; trade raises the
number of varieties and real wages; larger markets attract more than their share
of an industry.
\tested The \omterm{def:b3:scale-variety-firms:home-market}{home-market effect} has had mixed tests. On U.S. and Canadian
manufacturing, Head and Ries (2001) found it in the variation between industries but
not within them, and judged that most of the evidence favoured a model in which
each country's goods differ for reasons other than scale; Davis and Weinstein
rejected the simple version on OECD data, then found that market access with
\omterm{def:b3:ricardian-trade:comparative-advantage}{comparative advantage} does shape OECD manufacturing. % ledger: b3sv:hr.hme, b3sv:dw.hme
The gain from variety is measured, not tested: imported varieties quadrupled in the
United States from 1972 to 2001, worth 2.8~\% of GDP computed with the model's own
formula. % ledger: b3sv:bw.variety
\breaks Every firm exports, while most firms in every industry sell only at home.
\end{model}

\section{Exporters are few, large and productive}\label{sec:b3:scale-variety-firms:exporters}

The \omterm{def:b3:scale-variety-firms:melitz}{heterogeneous-firms model} predicts not only that few firms export but which
ones: the largest and most productive. The census data agree. Among United
States manufacturing firms in 2002, exporters had 1.19 log points more employment
than non-exporters --- about 3.3 times as many workers --- 1.48 log points more
shipments and 0.26 more value added per worker. Within the same industry the
differences shrink but remain: 0.97 log points of employment, 0.11 of value added
per worker, 0.06 of wages, all significant at the 1~\% level.
% ledger: b3sv:bjrs.t3
Part of the premium in value added per worker is capital and skill: exporters
use more of both. But the model asks the data a harder question: did exporters
become productive by exporting, or did productive firms start to export? Studies
that follow firms over time find that firms that start exporting were already
larger and more productive before they did (\cref{exo:b3:scale-variety-firms:9}).

\section{Reallocation after a trade agreement}\label{sec:b3:scale-variety-firms:reallocation}

\begin{record}[The Canada--U.S. free trade agreement, 1989]
Trefler studied Canadian manufacturing from 1989 to 1996, after an agreement that
cut tariffs without other reforms at the same time. In the industries with the
largest Canadian tariff cuts, employment fell by 15~\%, and output and the number
of plants by about 10~\%: the short-run cost. Labour productivity in the same
industries rose by 17~\% over the period, a compounded 2.1~\% a year, and about half of that rise came from
plant turnover --- entry and exit --- and rising technical efficiency.
% ledger: b3sv:tre.fta
\end{record}

Chile had opened its economy earlier, in the late 1970s and early 1980s.
Pavcnik's study of its manufacturing plants found productivity gains within
plants, especially in import-competing sectors, and that much of the aggregate
gain came from moving resources and output from less to more efficient
producers. % ledger: b3sv:pav.chile
Both records show the selection the model predicts, and both show its cost: the
plants that close and the jobs they take with them.

\begin{definition}[Foreign direct investment]\label{def:b3:scale-variety-firms:fdi}
\index{foreign direct investment}
\emph{Foreign direct investment} (FDI) is a firm's acquisition or creation of
production abroad that it controls. It is \emph{horizontal} when the foreign
affiliate makes the same goods to sell in its own market, and \emph{vertical}
when stages of production are placed abroad to use cheaper inputs, the output
being shipped on.
\end{definition}

A firm can serve a foreign market by exporting, paying the trade cost on every
unit, or by building an affiliate there, paying a larger fixed cost and no
trade cost. Extending the \omterm{def:b3:scale-variety-firms:melitz}{heterogeneous-firms model}, Helpman, Melitz and Yeaple
predicted that the least productive firms sell at home, the next export, and the
most productive set up affiliates. With United States exports and affiliate
sales in 38 countries and 52 sectors, they found that transport costs and tariffs
shift sales from exports to affiliates, and that sectors in which firms differ
more have more affiliate sales relative to exports. % ledger: b3sv:hmy.fdi
Offshoring, the vertical kind, cuts the chain of production into pieces and
places each where it is cheapest: a firm's ``imports'' are then often its own
parts coming home.

\begin{model}{Heterogeneous firms}
\assumes Firms draw productivity after paying a sunk entry cost; exporting has a
fixed cost as well as a trade cost; CES demand; free entry.
\predicts Only the most productive firms export; exporters are larger;
opening to trade drives the least productive firms out and raises average
industry productivity without any firm improving.
\tested Fewer than one United States manufacturing firm in five exports, and
exporters are larger and more productive; after the Canada--U.S. agreement,
labour productivity rose 17~\% in the most exposed industries, half of it through
plant turnover and efficiency; in Chile, liberalisation moved output to more
efficient plants. % ledger: b3sv:bjrs.t2.all, b3sv:bjrs.t3, b3sv:tre.fta, b3sv:pav.chile
\breaks Firms differ in quality and in the markets they sell to, not only in
productivity; within-plant gains, such as those found in Chile, are outside the
model.
\end{model}

\section{Exercises}\label{sec:b3:scale-variety-firms:exercises}

\begin{exercise}[$\star$]\label{exo:b3:scale-variety-firms:1}
In the opening chart, which industries have the largest and the smallest share
of firms that export? In which industries is the share above that of all
manufacturing? What do these industries have in common?
\end{exercise}

\begin{exercise}[$\star$]\label{exo:b3:scale-variety-firms:2}
Brenland's exports and imports with Kestria, in millions of dalens, are 40 and
30 for furniture, 10 and 50 for machines, 60 and 0 for grain. Compute the
\omterm{def:b3:scale-variety-firms:intra-industry}{Grubel--Lloyd index} for each industry and the aggregate index. Which trade would
Ricardo's model explain?
\end{exercise}

\begin{exercise}[$\star$]\label{exo:b3:scale-variety-firms:3}
With CES demand and $\sigma = 4$, find a firm's price as a multiple of marginal
cost, and its Lerner index. How do both change if varieties become closer
substitutes?
\end{exercise}

\begin{exercise}[$\star$]\label{exo:b3:scale-variety-firms:4}
Give an observation that would contradict the \omterm{def:b3:scale-variety-firms:home-market}{home-market effect}, and explain
why the \omterm{def:b3:factor-endowments:ho}{Heckscher--Ohlin model} predicts the opposite of the \omterm{def:b3:scale-variety-firms:home-market}{home-market effect}
for a good that consumers in one country especially like.
\end{exercise}

\begin{exercise}[$\star\star$]\label{exo:b3:scale-variety-firms:5}
A country has $L = 1{,}000$ workers. A firm needs $F = 10$ workers to start and $c =
1$ worker per unit; $\sigma = 5$ and the wage is 1. Find the price, output per firm,
employment per firm, the number of firms, and the real wage $1/P$.
\end{exercise}

\begin{exercise}[$\star\star$]\label{exo:b3:scale-variety-firms:6}
The country of \cref{exo:b3:scale-variety-firms:5} opens to free trade with an
identical country. Find the number of varieties consumed, the price and the
percentage gain in the real wage. Which is larger: the gain here, or with
$\sigma = 3$?
\end{exercise}

\begin{exercise}[$\star\star$]\label{exo:b3:scale-variety-firms:7}
With $\sigma = 3$, Pareto productivity of shape $k = 6$ (minimum 1), $f = 10$ and $f_e =
1$, find the closed-economy cut-off numerically, from a large sample of draws,
by setting an entrant's mean profit equal to $f_e$; compare it with the formula
of \cref{prop:b3:scale-variety-firms:cutoff}, and find the share of entrants that
produce.
\end{exercise}

\begin{exercise}[$\star\star$]\label{exo:b3:scale-variety-firms:8}
Productivity is Pareto with $k = 3$ and minimum 1. The domestic cut-off is 1.5 and
the \omterm{def:b3:scale-variety-firms:melitz}{export cut-off} 3. What share of entrants produce, and what share of
producers export? What would the export share be with $k = 2$?
\end{exercise}

\begin{exercise}[$\star\star$]\label{exo:b3:scale-variety-firms:9}
What data on firms would contradict selection into exporting? Describe a test,
and the alternative explanation it would support.
\end{exercise}

\begin{exercise}[$\star\star\star$]\label{exo:b3:scale-variety-firms:10}
Derive the CES demand \eqref{eq:b3:sv:demand} from \eqref{eq:b3:sv:ces}, show that
$U = E/P$, and deduce the constant markup of
\cref{prop:b3:scale-variety-firms:markup}.
\end{exercise}

\begin{exercise}[$\star\star\star$]\label{exo:b3:scale-variety-firms:11}
With $\sigma = 4$, $k = 5$, $f = 1$, $f_e = 0.5$, $\tau = 1.3$ and $f_x = 1.5$, simulate
100{,}000 entrants. Compute industry productivity --- the revenue-weighted mean of
the producers' $\varphi$ --- before and after opening, and split the rise into the
part due to exit and the part due to market share moving to exporters.
\end{exercise}

\begin{exercise}[$\star\star\star$]\label{exo:b3:scale-variety-firms:12}
In the economy of the listing ($f_x = 90$), the variable trade cost $\tau$ falls from
1.2 to 1.1. Find the new domestic cut-off and the share of producers that export
(the extensive margin), and the change in the export sales of a firm with $\varphi = 4$
(the intensive margin). Compare with a cut in $f_x$, in the weekend problem.
\end{exercise}

\section{Problem: Brenland's Furniture Makers}\label{sec:b3:scale-variety-firms:problem}

\begin{problem}[{Weekend problem --- Brenland's furniture makers}]\label{pb:b3:scale-variety-firms:furniture}
Brenland's furniture industry is a heterogeneous-firms industry with $\sigma = 3$.
Productivity $\varphi$, output per worker, is Pareto with shape $k = 4$ and minimum 1:
$\Pr(\varphi > z) = z^{-4}$. The wage is 1 (thousand dalens a year). Producing costs $f = 40$
a year and entry costs $f_e = 2.5$ a year (in thousands of dalens). Kestria is a
furniture market of the same size; selling there costs $\tau = 1.2$ and $f_x = 90$ a
year. % figdata: furn_params
An entrant's expected profit, when the domestic cut-off is $\varphi^*$ and there is no
trade, is $f\varphi^{*-k}(\sigma - 1)/(k - \sigma + 1)$.

\medskip\noindent\textbf{Part I --- the closed industry.}
\begin{enumerate}
\item Find the price of a firm with productivity $\varphi$, as a multiple of its
marginal cost $1/\varphi$.
\item Find the revenue of a firm at the domestic cut-off.
\item Revenue is proportional to $\varphi^{\sigma-1}$. Show that a producer's profit is
$40\bigl((\varphi/\varphi^*)^2 - 1\bigr)$.
\item Find the cut-off from free entry.
\item What share of entrants produce? Given that the mean of a Pareto variable
above $z$ is $kz/(k - 1)$, find the average productivity of producers.
\item Find the revenue and the profit of a firm with $\varphi = 4$.
\end{enumerate}

\medskip\noindent\textbf{Part II --- opening to Kestria.}
\begin{enumerate}[resume]
\item Show that the \omterm{def:b3:scale-variety-firms:melitz}{export cut-off} is $1.8$ times the domestic cut-off.
\item With trade, free entry gives $\varphi^{*4} = 16\bigl(1 + \tau^{-4}f/f_x\bigr)$. Find the
two cut-offs.
\item What share of the firms that produced in autarky exit?
\item What share of producers export? Compare with the furniture industry of the
opening chart.
\item By what factor are the total sales of a firm at the \omterm{def:b3:scale-variety-firms:melitz}{export cut-off} larger
than those of a firm at the domestic cut-off?
\item Find the domestic sales, exports and total sales of the firm with $\varphi =
4$, and compare with autarky.
\item What happens to the sales of a firm with $\varphi = 3$, and why?
\end{enumerate}

\medskip\noindent\textbf{Part III --- a trade agreement halves the fixed export cost.}
\begin{enumerate}[resume]
\item With $f_x = 45$, find the ratio $\varphi_x/\varphi^*$ and the two cut-offs.
\item What share of producers export now?
\item What share of the producers before the agreement exit?
\item What happens to the firm with $\varphi = 3$? Find its total sales.
\item Find the change in the exports of the firm with $\varphi = 4$. Explain the sign.
\item Do Brenland's total furniture exports rise? Through which margin?
\item By what percentage does the average productivity of producers rise?
\end{enumerate}

\medskip\noindent\textbf{Part IV --- productivity and welfare.}
\begin{enumerate}[resume]
\item No firm's productivity changed. Where does the rise in average productivity
come from?
\item In this model a worker's real income is proportional to $\varphi^*$. Find the
welfare gain from opening, from the agreement, and in all.
\item What would a Krugman model with identical firms predict for productivity and
for the source of the gains?
\item Which of these predictions does the Canada--U.S. record support?
\item Who loses from the agreement, and in what sense?
\item State the results: the rise in average productivity after the agreement,
and the share of producing firms that exit.
\end{enumerate}
\end{problem}
